\begin{abcliste}
\abc
\begin{align}
 \alpha &= \alF = \arcsin\frac{\la}{\bs} = \al \approx \result{\alP{0}{2}}
\end{align}

\abc A photon has the momentum $p = h/\lambda$, so
\begin{align}
 \Delta p_y &\approx \frac{h}{\lambda}\cdot\frac{\lambda}{b} = \pyF = \frac{\nch}{\bs} = \py \approx \result{\pyP{0}{2}}
\end{align}
With $\Delta y \approx b$: $\Delta y\cdot\Delta p_y \approx h$, which is more than $h/(4\pi) = \hbP{0}{2}$, as Heisenberg demands.

\abc Half the width: $\sin\alpha$ doubles, and the central maximum gets \result{\text{about twice as wide}}. The better the position sideways is fixed, the more the momentum sideways spreads.
\end{abcliste}
